\documentclass{article}

% if you need to pass options to natbib, use, e.g.:
%     \PassOptionsToPackage{numbers, compress}{natbib}
% before loading neurips_2024


% ready for submission
%\usepackage{neurips_2024}
\usepackage{graphicx}   % <--- 必须改为 graphicx
\usepackage{amsmath}
\usepackage{subcaption} % <--- 确保这一行也在，否则子图无法运行
\usepackage{algorithm}
\usepackage{algorithmic}


% to compile a preprint version, e.g., for submission to arXiv, add add the
% [preprint] option:
%\usepackage[preprint]{neurips_2024}


% to compile a camera-ready version, add the [final] option, e.g.:
     \usepackage[final]{neurips_2024}


% to avoid loading the natbib package, add option nonatbib:
%    \usepackage[nonatbib]{neurips_2024}


\usepackage[utf8]{inputenc} % allow utf-8 input
\usepackage[T1]{fontenc}    % use 8-bit T1 fonts
\usepackage{hyperref}       % hyperlinks
\usepackage{url}            % simple URL typesetting
\usepackage{booktabs}       % professional-quality tables
\usepackage{amsfonts}       % blackboard math symbols
\usepackage{nicefrac}       % compact symbols for 1/2, etc.
\usepackage{microtype}      % microtypography
\usepackage{xcolor}         % colors
\usepackage{subcaption}


\title{Visual Illusion Generation System Based on Diffusion Models}


% The \author macro works with any number of authors. There are two commands
% used to separate the names and addresses of multiple authors: \And and \AND.
%
% Using \And between authors leaves it to LaTeX to determine where to break the
% lines. Using \AND forces a line break at that point. So, if LaTeX puts 3 of 4
% authors names on the first line, and the last on the second line, try using
% \AND instead of \And before the third author name.


\author{%
  Zihan Yang \\
  EECS, Peking University \\
  \texttt{email} \\
  \And
  Yuming Fang \\
  EECS, Peking University \\
  \texttt{email} \\
  \And
  Yueran Wang \\
  SMS, Peking University \\
  \texttt{email} \\
}


\begin{document}


\maketitle


\begin{abstract}
  For this Computer Vision final project, we selected the topic of Creating Visual
  Cognitive Illusions for our experimentation.This paper presents an in-depth study on visual illusion generation using diffusion models. Building upon the reproduction and analysis of the existing Diffusion Illusion model, we identified its implementation flaws and theoretical limitations regarding complex geometric transformations and multi-stable perception. To address these issues, we constructed an Enhanced Illusion Generation Framework. By leveraging the generative prior of Score Distillation Sampling (SDS) and integrating a novel set of differentiable transformation operators into the optimization loop, our framework enables the decoupled optimization of image parameters across frequency domains, color spaces, and geometric perspectives. We successfully extended and implemented [7] distinct categories of visual illusions, including multi-scale hybrids, color-channel separation, cross-modal color-grayscale camouflage, and motion-induced parallax.To enhance reproducibility and interactivity, we developed a dedicated project homepage and an online execution environment. Experimental results demonstrate that our method significantly expands the flexibility and quality of generative illusions, providing new technical avenues for interactive art creation and multi-modal information hiding.
  
\end{abstract}


\section{Introduction}

Visual illusions—images designed to deceive the human visual system—encompass various categories, including ambiguous illusions, geometric-optical illusions, and motion illusions. Historically, they have remained a focal point of research in both psychology and computer graphics. In recent years, Diffusion Models have demonstrated exceptional capability in modeling complex distributions. Their scope of application has expanded far beyond mere image synthesis, extending to general robotic action trajectory generation\cite{chi2025diffusion} and Sim-to-Real navigation strategies (e.g., NavDP\cite{cai2025navdp}). However, despite the revolutionary breakthroughs Generative AI has brought to image creation, precise control over the geometric features and semantic attributes required for visual illusions remains a significant challenge.

Building upon the "Diffusion Illusion" framework by Ryan\cite{10.1145/3641519.3657500}, we identified theoretical limitations in existing models, specifically their restriction to a narrow range of illusion types and their reliance solely on text-based inputs. To address this, we focus on ambiguous illusions and extend the framework by incorporating an image-guidance mechanism and a suite of differentiable transformation operators. This approach not only enriches the diversity of generatable illusions but also supports user-supplied reference images. Specifically, we have implemented the following effects:

\textbf{Contributions.} Our main contributions are summarized as follows:

\begin{itemize}
    \item \textbf{Reproduction and Enhancement:} We successfully reproduced the three baseline illusions from the original \emph{Diffusion Illusions} framework and introduced critical improvements to their stability and visual quality, addressing limitations in the original implementation.
    
    \item \textbf{Automated Evaluation Pipeline:} We constructed a novel closed-loop pipeline that integrates automated generation with Vision-Language Model (VLM) based scoring. This system enables the objective quantification of illusion effectiveness and facilitates the high-throughput screening of optimal samples, replacing subjective manual selection.
    
    \item \textbf{Physical World Validation:} We bridged the gap between digital synthesis and physical perception by printing our generated samples. Experiments confirm that our illusions—particularly the cylindrical anamorphosis and hybrid images—maintain their perceptual efficacy in real-world environments under varying viewing conditions.
    
    \item \textbf{Framework Abstraction and Extension:} We formalized the previous specific implementations into a \textbf{Unified Multi-View Optimization Framework}. Based on this abstraction, we significantly expanded the solution space by introducing eight novel illusion categories, the implementation details of which are comprehensively described in Section \ref{sec:method}.

    \item  \textbf{}
\end{itemize}


\section{Related Work}
\label{related_work}

\subsection{Traditional Methods and Controlled Generation}

Early illusion generation relied heavily on frequency analysis or manual geometric design. Oliva\cite{10.1145/1141911.1141919}.'s Hybrid Images exploited multi-scale perception to blend frequency bands, while works on Camouflage and Anamorphosis \cite{chu2010camouflage} focused on texture optimization and geometric projection. With the rise of generative AI, methods like ControlNet\cite{zhao2023uni} and Illusion Diffusion\cite{10.1145/3641519.3657500} enabled the embedding of hidden patterns (e.g. QR codes) via spatial conditioning. However, these inference-based approaches often struggle to generate multi-stable illusions that require consistency across varying viewpoints or frequency domains.

\subsection{Optimization-based Paradigms}

To overcome single-view limitations, optimization-based frameworks leverage Score Distillation Sampling (SDS) to enforce multi-view consistency. Diffusion Illusions\cite{10.1145/3641519.3657500} pioneered optimizing a single image against prompts under affine transformations. Visual Anagrams\cite{geng2024visual} extended this using pixel-space models for rearrangement illusions, while DreamFusion\cite{poole2022dreamfusion} applied SDS to 3D generation. Our work builds on these foundations but introduces specialized decoupling operators for frequency and color spaces, addressing the limitations of current methods in handling multi-scale and cross-modal illusions.

\section{Data}
\label{data}


As our method is a zero-shot optimization approach, we do not require a traditional training dataset. Instead, we utilize:

Pre-trained Model: Stable Diffusion (based on the work of Diffusion Illusion\cite{10.1145/3641519.3657500}) acts as the source of visual knowledge and DeepFloyd IF.

Input Prompts: beyond the example prompts provided by the original paper’s authors, we carefully designed a series of prompts with striking contrasts for illusion generation, which yielded favorable results. In addition, we incorporated images generated by nanobanana as prompt inputs and embedded them into the illusions.


\section{Method}
\label{sec:method}

\subsection{Unified Multi-View Optimization Framework}

This study is deeply inspired by "Diffusion Illusions" proposed by Burgert\cite{10.1145/3641519.3657500} Through an in-depth analysis of the original code, we abstract and formalize the generation process of such visual illusions into a \textbf{Unified Multi-View Optimization Framework}.

We parameterize the object to be optimized as $\theta$ (which can be a single image, multi-layer textures, or a 3D voxel grid). We define a set of differentiable transformation functions $\mathcal{T} = \{T_1, \dots, T_K\}$, where each $T_k$ represents a specific physical observation modality or image processing operation. Our objective is to find a parameter set $\theta$ such that all transformed views $v_k = T_k(\theta)$ simultaneously satisfy their respective target constraints $y_k$.

The general pseudocode for this framework is presented in Algorithm~\ref{alg:framework}.

\begin{algorithm}[t]
\caption{Unified Framework for Diffusion-Based Visual Illusions}
\label{alg:framework}
\begin{algorithmic}[1]
% 在 REQUIRE 中增加了 total iterations T
\REQUIRE Pre-trained Diffusion Model $\epsilon_\phi$, Set of Transformations $\mathcal{T} = \{T_1, \dots, T_K\}$, Set of Targets $\mathcal{Y} = \{y_1, \dots, y_K\}$, Total iterations $T$
\STATE \textbf{Initialize:} Learnable parameters $\theta$ (e.g., Fourier features, voxel grid, or texture layers)
\STATE \textbf{Optimizer:} Initialize optimizer (e.g., Adam, SGD) for $\theta$
% --- 修改开始：将 While 改为 For ---
\FOR{$t = 1$ to $T$}
    \STATE $x \leftarrow \text{Construct}(\theta)$ \COMMENT{Synthesize base object from parameters}
    \STATE $\mathcal{L}_{\text{total}} \leftarrow 0$
    \FOR{$k = 1$ to $K$}
        \STATE $v_k \leftarrow T_k(x)$ \COMMENT{Apply differentiable transformation}
        \IF{$y_k$ is a Text Prompt}
            % 这里可以对应你代码中的 model_sd.train_step
            \STATE $\mathcal{L}_k \leftarrow \mathcal{L}_{\text{SDS}}(v_k, y_k, \epsilon_\phi)$ 
        \ELSIF{$y_k$ is a Target Image}
            % 这里对应你代码中的 loss_secret (MSE)
            \STATE $\mathcal{L}_k \leftarrow ||v_k - y_k||^2_2$ 
        \ENDIF
        \STATE $\mathcal{L}_{\text{total}} \leftarrow \mathcal{L}_{\text{total}} + \lambda_k \cdot \mathcal{L}_k$
    \ENDFOR
    \STATE $\theta \leftarrow \theta - \eta \cdot \nabla_\theta \mathcal{L}_{\text{total}}$ \COMMENT{Update parameters via backpropagation}
\ENDFOR
% --- 修改结束 ---
\RETURN Optimized object $x$
\end{algorithmic}
\end{algorithm}

\subsection{Baseline Implementation}
We first summarize the three baseline implementations by Burgert\cite{10.1145/3641519.3657500}, which constitute special cases of this framework:

\begin{itemize}
    \item \textbf{Flippy Illusion:} $\theta$ is a single image. The transformation $T$ is a $180^\circ$ rotation. The image conforms to two distinct text descriptions when viewed upright and inverted.
    \item \textbf{Rotation Overlays:} $\theta$ contains top and bottom image layers. The transformation $T$ is defined as the superposition product of the two layers, where the top layer is rotated by multiples of $90^\circ$.
    \item \textbf{Hidden Characters:} $\theta$ consists of four independent images. The transformation $T$ is the pixel-wise product of the four images, revealing hidden information upon superposition.
\end{itemize}

\subsection{Our Improvements}
Based on the unified framework described above, we expanded the boundaries of the transformation function $T(\cdot)$ and proposed eight novel visual illusion applications.

\begin{enumerate}
    \item \textbf{Image-Driven Hard Constraint Optimization:} Addressing the limitation of the original framework which relied solely on stochastic text prompts (via SDS Loss), we introduced a deterministic image-guidance mechanism. By incorporating a pixel-wise Mean Squared Error (MSE) loss into the optimization loop, we allow the system to use user-provided images as rigid targets. This enables precise control over the generated illusion, such as forcing the superposition of multiple layers to reconstruct a specific logo or reference image rather than a generic semantic concept (Figure \ref{fig:imp1}).

    \item \textbf{Multi-Angle Moire Cryptography:} Unlike the rigid 90-degree rotation in the baseline, we introduce arbitrary affine transformations combined with specific target MSE constraints. We optimize a dual-layer texture structure (Base Layer and Decoder Layer) such that when superimposed at specific angles ($0^\circ, 120^\circ, 240^\circ$), high-frequency texture interference (Moiré effect) accurately decodes distinct hidden visual symbols, achieving rotation-angle-based visual decryption (Figure \ref{fig:imp2}).

    \item \textbf{Differentiable Cylindrical Anamorphosis:} We design $T$ as an optical mapping function simulating cylindrical reflection. Using differentiable Grid Sample techniques, we map the planar canvas to the cylindrical mirror coordinate system. During optimization, the original canvas is constrained to exhibit abstract textures formed by stretching, while the reflected image after $T_{\text{mirror}}$ transformation is constrained to exhibit concrete target images. This achieves inverse rendering of physical optics (Figure \ref{fig:imp3}).

    \item \textbf{Motion Integration Steganography:} Leveraging the persistence of vision effect, we define $T$ as the directional motion blur convolution kernel. The optimization goal is to make the static image present a visible high-frequency texture, while the blurred image after rapid shaking (simulated convolution) reveals a hidden target. This effectively encodes target information into the image's noise patterns, which is decodable only under temporal integration (Figure \ref{fig:imp4}).

    \item \textbf{Orthogonal Voxel Projection Synthesis:} We elevate $\theta$ to a $64^3$ resolution 3D Voxel Grid. The transformation $T$ is defined as volume projection (Mean Projection) along the three orthogonal axes X, Y, and Z. This allows us to decompose the 3D generation problem into three 2D diffusion guidance problems, thereby generating an ``impossible'' 3D structure that presents completely different objects from different orthogonal perspectives (Figure \ref{fig:imp5}).

    \item \textbf{Distance-Dependent Spectral Hybridization:} Inspired by Hybrid Images theory, we explicitly decompose $\theta$ into multiple frequency layers. $T$ is a set of Gaussian filters with different $\sigma$ (low-pass, band-pass, high-pass). We separately optimize low-frequency (far view), mid-frequency (mid view), and high-frequency (near view) components to correspond to different semantic descriptions. This achieves continuous semantic transition based on viewing distance within a single image (Figure \ref{fig:imp6}).

    \item \textbf{Cross-Domain Luminance Decoupling:} Utilizing the orthogonality of color perception, we define $T$ as the luminance conversion formula from RGB to grayscale ($Y = 0.299R + 0.587G + 0.114B$). We simultaneously impose uncorrelated semantic constraints on the original color image and the transformed grayscale image. Through optimization, we generate visual adversarial images that present one object in color mode but abruptly mutate into another in desaturated (black and white) mode (Figure \ref{fig:imp7}).

    \item \textbf{Intra-Channel Frequency Splitting:} We refine the frequency separation strategy into the RGB channels. The transformation $T$ extracts high and low-frequency components for each of the R, G, and B channels. This allows us to encode non-interfering visual information separately within the macro color block distribution (low frequency) and micro detail textures (high frequency) of the same color image, achieving finer texture nesting than traditional hybrid images (Figure \ref{fig:imp8}).
\end{enumerate}

\begin{figure}[t]
    \centering
    % --- 第一排 (Row 1): 3 Images ---
    \begin{subfigure}[b]{0.32\linewidth}
        \centering
        \includegraphics[width=\linewidth]{figures/1for8.pdf} 
        \caption{Hidden QR Code}
        \label{fig:imp1}
    \end{subfigure}
    \hfill
    \begin{subfigure}[b]{0.32\linewidth}
        \centering
        \includegraphics[width=\linewidth]{figures/2for8.pdf}
        \caption{RGB Moire Cryptography}
        \label{fig:imp2}
    \end{subfigure}
    \hfill
    \begin{subfigure}[b]{0.32\linewidth}
        \centering
        \includegraphics[width=\linewidth]{figures/3for8.pdf}
        \caption{Cylindrical Anamorphosis}
        \label{fig:imp3}
    \end{subfigure}
    
    \vspace{1em} % 行间距

    % --- 第二排 (Row 2): 3 Images ---
    \begin{subfigure}[b]{0.32\linewidth}
        \centering
        \includegraphics[width=\linewidth]{figures/4for8.pdf}
        \caption{Motion Steganography}
        \label{fig:imp4}
    \end{subfigure}
    \hfill
    \begin{subfigure}[b]{0.32\linewidth}
        \centering
        \includegraphics[width=\linewidth]{figures/5for8.pdf}
        \caption{Orthogonal 3D Projection}
        \label{fig:imp5}
    \end{subfigure}
    \hfill
    \begin{subfigure}[b]{0.32\linewidth}
        \centering
        \includegraphics[width=\linewidth]{figures/6for8.pdf}
        \caption{Hybrid Expression}
        \label{fig:imp6}
    \end{subfigure}

    \vspace{1em} % 行间距

    % --- 第三排 (Row 3): 2 Images (Centered) ---
    \begin{subfigure}[b]{0.32\linewidth}
        \centering
        \includegraphics[width=\linewidth]{figures/7for8.pdf}
        \caption{Color-Grayscale Switch}
        \label{fig:imp7}
    \end{subfigure}
    \hspace{0.02\linewidth} % 两图之间的微小间距
    \begin{subfigure}[b]{0.32\linewidth}
        \centering
        \includegraphics[width=\linewidth]{figures/8for8.pdf}
        \caption{Color Temperature Shift}
        \label{fig:imp8}
    \end{subfigure}

    \caption{\textbf{Gallery of Generated Illusions.} We demonstrate eight distinct visual illusion categories enabled by our unified framework. 
    (a) Four images (Miku, Frog, Cat, Giraffe) superimposed to form a functional QR code for our project homepage. 
    (b) Two layers reveal a red "PKU", green "EECS", and blue "SMS" logo when superimposed at $0^\circ$, $120^\circ$, and $240^\circ$ respectively. 
    (c) (Left) The planar "starry sky" canvas. (Center) The reflection visible on the cylindrical mirror. (Right) The target image (movie poster from \textit{Your Name}). 
    (d) (Left) Image of buildings. (Center) The effect after simulated shaking (motion blur). (Right) The target image (a rose). 
    (e) A generated 3D asset appearing as a butterfly, a tree, and an apple from three orthogonal viewing directions. 
    (f) (Left) A non-smiling person. (Right) The downscaled version appears to be smiling. 
    (g) (Left) A colorful landscape with a river and mountains. (Right) The grayscale version reveals a tiger. 
    (h) (Left) The image shows cool color tones. (Right) The downscaled version shifts to warm color tones.}
    \label{fig:all_methods}
\end{figure}

\section{Experiments}

\subsection{Experimental Settings}
To evaluate the quality of visual illusions generated by Diffusion Models, one viable approach is to employ Vision Language Models (VLMs) to verify whether the generated illusion images convey the intended semantics. Several studies have explored this direction, such as \emph{IllusionVQA} \cite{shahgir2024illusionvqa} and \emph{GVIL} \cite{zhang2023grounding}.

In this assignment, we designed an automated pipeline for the generation and evaluation of visual illusions. The framework of this pipeline consists of two main stages: automated generation using Diffusion Models and intelligent evaluation using VLMs. Under this framework, we evaluated and analyzed two types of illusions:
\begin{enumerate}
    \item \textbf{Text Flip Illusion (Ambigrams):} An image that reveals one string of text when viewed upright and a different string when rotated by 180 degrees \cite{geng2024visual}.
    \item \textbf{Hybrid Illusion:} As described in Section \ref{sec:methods}, where perception changes based on frequency or color modality.
\end{enumerate}

\subsection{Pipeline Architecture and Implementation}
Our experimental pipeline operates in two sequential stages, each equipped with dedicated scripts and producing structured output files. All scripts, experimental results (including generated images), natural language evaluations from the VLM, and numerical scores (stored as JSON files) have been uploaded to our project's GitHub repository.

\textbf{Stage 1: Automated Generation.}
For the Text Flip Illusion, we utilized DeepFloyd IF as the backbone model due to its superior character generation capabilities. For the Hybrid/Blurry Illusions, we employed Stable-Diffusion-v1-4. The generation process yields a structured output directory containing all candidate image samples.

\textbf{Stage 2: VLM-Based Intelligent Evaluation.}
Each generated candidate image was evaluated using \textbf{Qwen3-VL-4B-Instruct}, a Vision Language Model capable of understanding both visual content and text. To ensure objectivity, we implemented a "blind test" strategy where the VLM performs recognition without prior knowledge of the target words.

\subsection{Experimental Results}
For the Text Flip Illusion, we conducted experiments on two word pairs: (EECS, SMS) and (ICS, LOVE). We generated 3,333 candidate images for the (EECS, SMS) pair and 1,923 candidates for the (ICS, LOVE) pair. Additionally, for the Hybrid Illusion, we prepared a pool of 20 pre-defined prompts; for each generation, a prompt was randomly selected from this pool, resulting in a total of 100 sets of generated images. Due to space constraints, we present only the most representative experimental results below.

Table \ref{tab:score_dist} illustrates the score distribution for the (EECS, SMS) pair. As observed, generating high-quality visual illusions of this type presents significant difficulty, with only a very small fraction of samples aligning well with the prompt. Figure \ref{fig:samples} (below) displays three sets of generated images which received VLM scores of 0, 5, and 9 respectively, clearly demonstrating the disparity in generation quality.

\begin{table}[h]
  \caption{VLM Score Distribution for (EECS, SMS) Illusion Generation}
  \label{tab:score_dist}
  \centering
  % 使用 small 字体确保宽度适应，同时设置两组数据列
  \small
  \begin{tabular}{ccc @{\hskip 1cm} ccc}
    \toprule
    % 表头：左边一组，右边一组
    VLM Score & Count & Percentage & VLM Score & Count & Percentage \\
    \cmidrule(r){1-3} \cmidrule(l){4-6}
    
    % 数据行：手动将原来的19行数据拆分为左右两列
    0.0 & 1,171 & 35.1\% & 5.0 & 68 & 2.0\% \\
    0.5 & 36    & 1.1\%  & 5.5 & 15 & 0.5\% \\
    1.0 & 143   & 4.3\%  & 6.0 & 106& 3.2\% \\
    1.5 & 50    & 1.5\%  & 6.5 & 28 & 0.8\% \\
    2.0 & 77    & 2.3\%  & 7.0 & 193& 5.8\% \\
    2.5 & 26    & 0.8\%  & 7.5 & 666& 20.0\% \\
    3.0 & 75    & 2.3\%  & 8.0 & 462& 13.9\% \\
    3.5 & 15    & 0.5\%  & 8.5 & 181& 5.4\% \\
    4.0 & 58    & 1.7\%  & 9.0 & 4  & 0.1\% \\
    4.5 & 13    & 0.4\%  &     &    &       \\ % 最后一行右边留空
    \bottomrule
  \end{tabular}
\end{table}

\begin{figure}[h]
  \centering
  %--- 子图 1 (Score 0) ---
  \begin{subfigure}[b]{0.32\linewidth}
    \centering
    % 注意：将[参数]和{文件名}放在同一行，防止报错
    \includegraphics[width=\linewidth]{figures/sample_256_views_bad.png}
    \caption{Score 0 Sample}
    \label{fig:score0}
  \end{subfigure}
  \hfill % 在子图之间添加弹性间距
  %--- 子图 2 (Score 5) ---
  \begin{subfigure}[b]{0.32\linewidth}
    \centering
    
\includegraphics[width=\linewidth]{figures/sample_256_views_medium.png}
    \caption{Score 5 Sample}
    \label{fig:score5}
  \end{subfigure}
  \hfill % 在子图之间添加弹性间距
  %--- 子图 3 (Score 9) ---
  \begin{subfigure}[b]{0.32\linewidth}
    \centering
    \includegraphics[width=\linewidth]{figures/sample_256_views_good.png}
    \caption{Score 9 Sample}
    \label{fig:score9}
  \end{subfigure}

  \caption{Comparison of generated samples with VLM scores. (a), (b), and (c) show examples with scores of 0, 5, and 9 respectively. Higher scores indicate clearer text legibility and better illusion consistency.}
  \label{fig:samples}
\end{figure}

\subsection{Discussion}
The analysis above indicates that our pipeline successfully provides an automated method for filtering high-quality visual illusions, effectively validating the scaling law in the context of generative search—where increasing the quantity of candidate generations significantly increases the probability of discovering high-quality outliers.

Furthermore, we offer some insights regarding future directions. A primary challenge in many generative tasks lies in designing the appropriate Loss Function. We hypothesize that VLM automated scoring could potentially serve as a direct Loss Function (i.e., a "Loss Prompt"). However, the current limitation is that such VLM feedback is non-differentiable. If this process could be made differentiable, we believe it would unlock significant applications in optimizing complex semantic constraints.

\section{Limitations}

While this study significantly expands the boundaries of diffusion-based visual illusion generation through strategies such as image-driven hard constraints, frequency separation, and cross-domain decoupling, our experiments reveal substantial limitations regarding model architecture, theoretical constraints, and optimization stability.

\subsection{Artifacts in Latent Space Transformations}
Our framework utilizes Stable Diffusion (SD) as the generative backbone, which relies on a Variational Autoencoder (VAE) to compress images into latent space. For tasks involving geometric transformations—such as Orthogonal Voxel Projection or Multi-Angle Moire Cryptography—manipulating the latent grid directly rearranges the spatial configuration of latent patches but fails to rotate the encoded content within the patches themselves. This discrepancy results in "thatching artifacts" (jagged edges at patch boundaries) upon decoding, disrupting the visual continuity of the illusion. Pixel-based architectures, such as DeepFloyd IF, demonstrate superior adaptability to such transformations but entail significantly higher computational costs.

\subsection{Gradient Conflicts in Projective Transformations}

In exploring Cross-Domain Luminance Decoupling (Color-Grayscale Switch), we encountered challenges more severe than those in traditional geometric illusions. Previous works like Visual Anagrams\cite{geng2024visual} rely primarily on orthogonal transformations (e.g.rotation), which preserve invertibility and noise statistics. In contrast, the conversion from RGB to grayscale is a non-invertible projective transformation ($\mathbb{R}^3 \to \mathbb{R}^1$). Our experiments indicate that simple channel splicing leads to semantic incoherence, while gradient-based joint optimization suffers from "structural competition." Since the luminance channel carries the majority of structural information, gradients from the color objective and the grayscale objective often conflict violently. This forces the optimization to sacrifice the naturalness of the color image to satisfy the grayscale constraint, frequently resulting in oversaturated artifacts or unnatural textures.

\subsection{Trade-off between Pixel Alignment and Semantic Fidelity}

Generating visual illusions essentially involves identifying the intersection of multiple semantic concepts on the image manifold. This is particularly challenging when introducing Image-Driven Hard Constraint Optimization, where we enforce strict pixel-level alignment with reference targets (e.g.specific logos). Such rigid mathematical constraints often force the generative process to deviate from the natural image distribution. Consequently, the optimization tends to compromise texture realism for geometric precision, yielding images that satisfy structural requirements but exhibit perceptual anomalies—such as high-frequency noise or over-sharpening—resembling adversarial examples.

\subsection{Over-optimization and Ghosting}

Contrary to the expectation that more iterations yield better results, we observed that image quality does not monotonically improve with iteration count in SDS-based optimization. Empirical evidence from our Distance-Dependent Spectral Hybridization experiments suggests that perceptual quality peaks at approximately 1500–2000 iterations. Extending optimization beyond this point (e.g., to the 5000 iterations used in baseline works) leads to "over-optimization," where the model introduces excessive high-frequency details to maximize classifier scores. This results in oversaturation and "ghosting" (where hidden views bleed into the visible view), indicating a lack of a perception-aware stopping mechanism in current algorithms.

\subsection{Hyperparameter Sensitivity}

By introducing complex transformations in both frequency and spatial domains—such as Intra-Channel Frequency Splitting and Motion Integration Steganography—our method exhibits high sensitivity to hyperparameters. Unlike standard text-to-image generation, illusion synthesis requires finding a Pareto optimal solution among competing loss functions. Critical parameters include spectral cutoff frequencies, loss weights for different views, and the optimization schedule. Experiments show that these parameters often require fine-grained, case-specific tuning to prevent one visual component from dominating the others, hindering fully automated generation.

\subsection{Dataset Bias and Generative Boundaries}

Our study finds that the generative capacity of diffusion models is significantly constrained by training data priors and encoder architecture, manifesting in three key aspects. First, since SD v1.4 was primarily trained on the LAION-5B dataset, the model exhibits significant Western-centric bias. This leads to severe semantic deviations when processing specific cultural symbols or complex non-Western scenarios, as well as inevitable Western stylistic biases in generated images. Second, the model has internalized strong 3D Consistency Priors of the physical world. When attempting to generate "geometric paradoxes" like Penrose triangles\cite{dodik2025meschers}, the model tends to use prior knowledge to "correct" spatial contradictions, forcing them into ordinary objects that conform to perspective laws, causing the geometric illusion to fail. Finally, limited by the CLIP text encoder's semantic rather than character-aware mechanism, the model shows instability in handling Typography. When generating text-based illusions such as Ambigrams, the model often treats text as texture rather than symbol sequences, frequently resulting in missing strokes or gibberish, making precise glyph control difficult to achieve.

\section{Conclusion}

In this project, we established a unified framework for visual illusion generation based on Score Distillation Sampling (SDS), significantly deepening our understanding of controllable generation within diffusion models. By injecting a novel set of differentiable transformation operators into the SDS optimization pipeline, we transcended the limitations of simple 2D rotations. Our work successfully leveraged the nuances of human perception—specifically the interplay between frequency, luminance, and chrominance—to realize ? distinct categories of illusions:

While achieving these results, we also identified a fundamental limitation of Stable Diffusion in illusion generation: its reliance on a Latent Diffusion Architecture. Since the latent code represents a compressed encoding of image "patches" rather than individual pixels, performing rigid transformations (like rotation) on the latent map disrupts the spatial structure required by the decoder. This lack of rotation equivariance in the latent space often leads to high-frequency artifacts and "ghosting" in multi-view illusions. Future work could address this by adopting pixel-space diffusion models (e.g. DeepFloyd IF) or integrating Visual Language Models (VLMs) for automated illusion evaluation. Due to space constraints, a more comprehensive analysis of the limitations and technical constraints is also included in the Appendix.

To facilitate reproducibility and interaction, we developed a project homepage and a Google Colab environment, allowing users to reproduce all generated illusions. Detailed documentation is provided in the Appendix.

\bibliographystyle{plain}
\bibliography{references}
\end{document}